\documentclass[../../main.tex]{subfiles}
\graphicspath{{\subfix{../../image/}}}

\begin{document}

\section{习题}\label{section:ch08-7-习题}

在下列习题中，$U$ 总表示 $\mathbb{R}^n$ 的一个有界开子集，边界光滑。除非另作说明，所有给定函数均假设光滑。

\begin{exercise}\label{exercise:ch08-7-1}
本问题说明一个弱收敛序列可能行为很坏。
\begin{enumerate}[(a)]
\item 证明 $u_k(x)=\sin(kx)\rightharpoonup 0$（当 $k\to\infty$）于 $L^2(0,1)$。
\item 固定 $a,b\in\mathbb{R}$，$0<\lambda<1$。定义
\begin{align*}
u_k(x)\triangleq\begin{cases}a&\text{若 }j/k\leqslant x<(j+\lambda)/k,\\ b&\text{若 }(j+\lambda)/k\leqslant x<(j+1)/k,\end{cases}\qquad (j=0,\cdots,k-1).
\end{align*}
证明 $u_k\rightharpoonup\lambda a+(1-\lambda)b$ 于 $L^2(0,1)$。
\end{enumerate}
\end{exercise}

\begin{exercise}\label{exercise:ch08-7-2}
设 $\phi:U\to\mathbb{R}$。求 $L=L(p,z,x)$，使 PDE
\begin{align*}
-\Delta u+D\phi\cdot Du=f\quad\text{在 }U\text{ 中}
\end{align*}
是与泛函 $I[w]\triangleq\int_U L(Dw,w,x)\,\mathrm{d}x$ 相关联的 Euler--Lagrange 方程。（提示：寻找一个带指数项的 Lagrangian。）
\end{exercise}

\begin{exercise}\label{exercise:ch08-7-3}
热方程的椭圆正则化是 PDE
\begin{align*}
u_t-\Delta u-\varepsilon u_{tt}=0\quad\text{在 }U_T\text{ 中},
\end{align*}
其中 $\varepsilon>0$，$U_T=U\times(0,T]$。证明上述方程是与某个能量泛函 $I_\varepsilon[w]\triangleq\iint_{U_T}L_\varepsilon(Dw,w_t,w,x,t)\,\mathrm{d}x\,\mathrm{d}t$ 相关联的 Euler--Lagrange 方程。（提示：寻找一个带涉及 $t$ 的指数项的 Lagrangian。）
\end{exercise}

\begin{exercise}\label{exercise:ch08-7-4}
设 $\eta:\mathbb{R}^n\to\mathbb{R}$ 为 $C^1$。
\begin{enumerate}[(a)]
\item 证明 $L(P,z,x)=\eta(z)\det P$（$P\in\mathbb{M}^{n\times n}$，$z\in\mathbb{R}^n$）是零 Lagrangian。
\item 推得：若 $\mathbf{u}:\mathbb{R}^n\to\mathbb{R}^n$ 为 $C^2$，则
\begin{align*}
\int_U\eta(\mathbf{u})\det D\mathbf{u}\,\mathrm{d}x
\end{align*}
只依赖于 $\mathbf{u}|_{\partial U}$。
\end{enumerate}
\end{exercise}

\begin{exercise}[续上题]\label{exercise:ch08-7-5}固定 $x_0\notin\mathbf{u}(\partial U)$，并选取如上函数 $\eta$ 使 $\int_{\mathbb{R}^n}\eta\,\mathrm{d}z=1$，$\operatorname{spt}\eta\subset B(x_0,r)$，$r$ 取充分小使 $B(x_0,r)\cap\mathbf{u}(\partial U)=\emptyset$。定义
\begin{align*}
\deg(\mathbf{u},x_0)=\int_U\eta(\mathbf{u})\det D\mathbf{u}\,\mathrm{d}x,
\end{align*}
即 $\mathbf{u}$ 关于 $x_0$ 的\textbf{度数}。证明该度数是整数。
\end{exercise}

\begin{exercise}\label{exercise:ch08-7-6}
设 $\Sigma\subset\mathbb{R}^3$ 表示光滑函数 $u:U\to\mathbb{R}$（$U\subset\mathbb{R}^2$）的图像。则
\begin{align*}
\int_U(1+|Du|^2)^{-3/2}\det D^2u\,\mathrm{d}x
\end{align*}
表示 Gauss 曲率在 $\Sigma$ 上的积分。证明此表达式只依赖于 $Du$ 在 $\partial U$ 上的限制。（微分几何中的 Gauss--Bonnet 定理以 $\partial\Sigma$ 的测地曲率计算上式。）
\end{exercise}

\begin{exercise}\label{exercise:ch08-7-7}
设 $m=n$。证明
\begin{align*}
L(P)=\operatorname{tr}(P^2)-\operatorname{tr}(P)^2\qquad (P\in\mathbb{M}^{n\times n})
\end{align*}
是零 Lagrangian。
\end{exercise}

\begin{exercise}\label{exercise:ch08-7-8}
解释为何 \S8.2 中的方法不能用于证明泛函
\begin{align*}
I[w]\triangleq\int_U(1+|Dw|^2)^{1/2}\,\mathrm{d}x
\end{align*}
在 $\mathcal{A}\triangleq\{w\in W^{1,q}(U)\mid w=g\text{ 在 }\partial U\text{ 上}\}$ 上的极小元存在性，对任何 $1\leqslant q<\infty$。
\end{exercise}

\begin{exercise}[方程组的二阶变分]\label{exercise:ch08-7-9}设 $\mathbf{u}:U\to\mathbb{R}^m$ 是泛函
\begin{align*}
I[\mathbf{w}]\triangleq\int_U L(D\mathbf{w},\mathbf{w},x)\,\mathrm{d}x
\end{align*}
的光滑极小元。
\begin{enumerate}[(a)]
\item 证明
\begin{align*}
\sum_{i,j=1}^n\sum_{k,l=1}^m\frac{\partial^2 L}{\partial p_i^k\partial p_j^l}(D\mathbf{u},\mathbf{u},x)\eta^k\eta^l\xi_i\xi_j\geqslant 0
\end{align*}
对所有 $x\in U$、$\xi\in\mathbb{R}^n$、$\eta\in\mathbb{R}^m$ 成立。
\item 给出一个非凸函数 $L:\mathbb{M}^{m\times n}\to\mathbb{R}$ 的例子，满足
\begin{align*}
\sum_{i,j=1}^n\sum_{k,l=1}^m\frac{\partial^2 L(P)}{\partial p_i^k\partial p_j^l}\eta^k\eta^l\xi_i\xi_j\geqslant 0
\end{align*}
对所有 $P\in\mathbb{M}^{m\times n}$、$\xi\in\mathbb{R}^n$、$\eta\in\mathbb{R}^m$ 成立。
\end{enumerate}
\end{exercise}

\begin{exercise}\label{exercise:ch08-7-10}
用 \S8.4.1 的方法证明：存在非平凡弱解 $u\in H_0^1(U)$，$u\not\equiv 0$，满足
\begin{align*}
\begin{cases}
-\Delta u=|u|^{q-1}u&\text{在 }U\text{ 中},\\
u=0&\text{在 }\partial U\text{ 上},
\end{cases}
\end{align*}
其中 $1<q<\frac{n+2}{n-2}$，$n>2$。
\end{exercise}

\begin{exercise}\label{exercise:ch08-7-11}
设 $\beta:\mathbb{R}\to\mathbb{R}$ 光滑，且
\begin{align*}
0<a\leqslant\beta'(z)\leqslant b\qquad (z\in\mathbb{R})
\end{align*}
（$a,b$ 为常数）。设 $f\in L^2(U)$。阐述 $u\in H^1(U)$ 是非线性边值问题
\begin{align*}
\begin{cases}
-\Delta u=f&\text{在 }U\text{ 中},\\
\frac{\partial u}{\partial\nu}+\beta(u)=0&\text{在 }\partial U\text{ 上}
\end{cases}
\end{align*}
的弱解意味着什么。证明存在唯一弱解。
\end{exercise}

\begin{exercise}\label{exercise:ch08-7-12}
设 $u$ 是面积积分
\begin{align*}
I[w]=\int_U(1+|Dw|^2)^{1/2}\,\mathrm{d}x
\end{align*}
在给定边界条件 $w=g$ 在 $\partial U$ 上及约束 $J[w]=\int_U w\,\mathrm{d}x=1$ 下的光滑极小元。证明 $u$ 的图像是一个常平均曲率曲面。（提示：回顾 \S8.1.2 例 4。）
\end{exercise}

\begin{exercise}\label{exercise:ch08-7-13}
设 $f\in L^2(U)$。证明对偶变分原理：
\begin{align*}
\min_{w\in H_0^1(U)}\int_U\frac12|Dw|^2-fw\,\mathrm{d}x=\max_{\substack{\boldsymbol{\xi}\in L^2(U;\mathbb{R}^n)\\\operatorname{div}\boldsymbol{\xi}=f}}-\frac12\int_U|\boldsymbol{\xi}|^2\,\mathrm{d}x.
\end{align*}
\end{exercise}

\begin{exercise}[多值 PDE]\label{exercise:ch08-7-14}证明 \S8.4.2 中障碍问题的变分不等式\eqref{align:ch08-4-eq26} 可改写为
\begin{align*}
-\Delta u+\beta(u-h)\ni f
\end{align*}
其中多值函数
\begin{align*}
\beta(z)\triangleq\begin{cases}0&\text{若 }z>0,\\(-\infty,0]&\text{若 }z=0,\\\emptyset&\text{若 }z<0.\end{cases}
\end{align*}
（另见第 9 章习题 3。）
\end{exercise}

\begin{exercise}[逐点梯度约束]\label{exercise:ch08-7-15}
\begin{enumerate}[(a)]
\item 证明存在唯一极小元 $u\in\mathcal{A}$，使
\begin{align*}
I[w]\triangleq\int_U\frac12|Dw|^2-fw\,\mathrm{d}x,
\end{align*}
其中 $f\in L^2(U)$，且
\begin{align*}
\mathcal{A}\triangleq\{w\in H_0^1(U)\mid |Dw|\leqslant 1\text{ a.e.}\}.
\end{align*}
\item 证明
\begin{align*}
\int_U Du\cdot D(w-u)\,\mathrm{d}x\geqslant\int_U f(w-u)\,\mathrm{d}x
\end{align*}
对所有 $w\in\mathcal{A}$ 成立。
\end{enumerate}
\end{exercise}

\begin{exercise}\label{exercise:ch08-7-16}
设 $n\geqslant 3$，$U$ 为包含 $0$ 的有界开集。证明 $\mathbf{u}\triangleq\frac{x}{|x|}$ 属于 $H^1(U;\mathbb{R}^n)$，且是映到球面 $S^{n-1}$ 的调和映射。即证明 $\mathbf{u}$ 是
\begin{align*}
\begin{cases}
-\Delta\mathbf{u}=|D\mathbf{u}|^2\mathbf{u}&\text{在 }U\text{ 中},\\
|\mathbf{u}|=1
\end{cases}
\end{align*}
的弱解。
\end{exercise}

\begin{exercise}\label{exercise:ch08-7-17}
设 $u,\widehat{u}\in H_0^1(U)$ 均为 Dirichlet 能量
\begin{align*}
I[w]\triangleq\int_U|Dw|^2\,\mathrm{d}x
\end{align*}
在约束 $\int_U w^2\,\mathrm{d}x=1$ 下的极小元。又设 $u,\widehat{u}>0$ 在 $U$ 内。沿提示给出 $u\equiv\widehat{u}$ 在 $U$ 中的新证明。
（提示：定义 $w=\left(\frac{u^2+\widehat{u}^2}{2}\right)^{1/2}$，$s=\frac{u^2}{u^2+\widehat{u}^2}$，$\eta=\frac{u^2+\widehat{u}^2}{2}$；并证明
\begin{align*}
|Dw|^2=\eta\left|\left|s\frac{Du}{u}+(1-s)\frac{D\widehat{u}}{\widehat{u}}\right|\right|^2.
\end{align*}
推得
\begin{align*}
|Dw|^2\leqslant\eta\left(s\left|\left|\frac{Du}{u}\right|\right|^2+(1-s)\left|\left|\frac{D\widehat{u}}{\widehat{u}}\right|\right|^2\right)=\frac12|Du|^2+\frac12|D\widehat{u}|^2,
\end{align*}
从而 $\frac{Du}{u}=\frac{D\widehat{u}}{\widehat{u}}$ a.e.。）（Belloni--Kawohl, Manuscripta Math.\ 109 (2002), 229--231）
\end{exercise}

\begin{exercise}\label{exercise:ch08-7-18}
设 $a_1,a_2$ 为 $\overline{U}$ 上光滑正函数，$a_1\leqslant a_2$。设 $u_1,u_2$ 为
\begin{align*}
\operatorname{div}(a_1Du_1)=0,\qquad \operatorname{div}(a_2Du_2)=0\quad\text{在 }U\text{ 中}
\end{align*}
的光滑解，且 $Du_2\neq 0$ a.e.。最后设
\begin{align*}
u_1=u_2,\quad \frac{\partial u_1}{\partial\nu}=\frac{\partial u_2}{\partial\nu},\quad a_1=a_2\quad\text{在 }\partial U\text{ 上}.
\end{align*}
证明在 $U$ 内 $a_1\equiv a_2$，$u_1\equiv u_2$。（提示：观察 $\int_U a_1|Du_1|^2\,\mathrm{d}x=\int_U a_2|Du_2|^2\,\mathrm{d}x$。）
\end{exercise}

\begin{exercise}[动量守恒]\label{exercise:ch08-7-19}给定非线性波方程 $\Box u+f(u)=0$ 的解 $u$，对变换 $\mathbf{x}(x,t,\tau)=(x+\tau e_k,t)$、$w(x,t,\tau)=u(x+\tau e_k,t)$ 应用 Noether 定理，计算动量密度 $p_k$ 与动量流 $\mathbf{j}_k$，使其满足守恒律
\begin{align*}
(p_k)_t-\operatorname{div}(\mathbf{j}_k)=0\qquad (k=1,\cdots,n).
\end{align*}
\end{exercise}

\begin{exercise}\label{exercise:ch08-7-20}
设 $u$ 在某区域 $U\subseteq\mathbb{R}^n$ 内调和，且 $B(0,R)\subset U$，$u(0)=0$，$u\not\equiv 0$。对 $0<r<R$ 定义函数
\begin{align*}
a(r)\triangleq\frac1{r^{n-1}}\int_{\partial B(0,r)}u^2\,\mathrm{d}S,\qquad b(r)\triangleq\frac1{r^{n-2}}\int_{B(0,r)}|Du|^2\,\mathrm{d}x.
\end{align*}
我们在 \S8.6.2 导出了单调性公式
\begin{align*}
b'=\frac{2}{r^{n-2}}\int_{\partial B(0,r)}u_r^2\,\mathrm{d}S.
\end{align*}
\begin{enumerate}[(a)]
\item 证明
\begin{align*}
a'=\frac{2}{r^{n-1}}\int_{\partial B(0,r)}uu_r\,\mathrm{d}S=\frac2r b.
\end{align*}
\item 证明
\begin{align*}
b^2\leqslant\frac{r}{2}ab'.
\end{align*}
\item 定义频率函数 $f\triangleq\frac b a$，导出 Almgren 单调性公式：$f'\geqslant 0$。
\item 接着证明 $\frac{a'}a\leqslant\frac\beta r$，从而
\begin{align*}
a(r)\geqslant\gamma r^\beta\qquad (0<r<R),
\end{align*}
其中 $\beta\triangleq\frac{2b(R)}{a(R)}$，$\gamma\triangleq\frac{a(R)}{R^\beta}$。这是对一个非常数调和函数在其取零点附近增长速度的下界估计。
\end{enumerate}
\end{exercise}

\end{document}
